\documentclass{article}
\usepackage[utf8]{inputenc}
\usepackage[T1]{fontenc}
\usepackage[margin=2cm]{geometry}
\usepackage{hyperref}
\sloppy
\begin{document}
\section*{Historical source}
Encyclopaedia Britannica scan, source row 2131. Parse 3.5 extracted the text. Super 3.5 VL corrected it from an enlarged crop of each block. It can still contain recognition errors.
Page headers, footers, and page numbers are removed.

\subsection*{Untitled continuation (E0182)}
\textit{Continues an entry from an earlier page.}

Bombay. AUGUST. SEPTEMBER.
Days of Quantity of Days of Quantity of
the month. rain in the month. rain in
In. Tenths. In. Tenths.
19 0 2 22 0 6
20 0 2 23 0 3
21 0 1 24 0 3
22 0 4 25 0 2
25 0 4 26 0 2
26 0 7 27-30 0 0
27 0 5 OCTOBER.
28 0 2 2-7 2 1
30 0 1 9 0
31 0 2 15 0 3
SEPTEMBER. 16, 17 0 1
1 0 2 13 The air fine and
3 0 3
clear, without a cloud, the
9 0 1
land and sea-breezes re-
9 4 gular.
10 0 5 While quantity of rain in
11 0 9
11 0 0 9
12 0 2 May 1
13 0 5 June 44 7
16 0 4 July 19 9
17 0 0 4 August 19 9
18 0 5 September 11 2
19 9 6 October 4 5
20 3 2
21 0 4 Total 110 3
In this journal our author makes no mention of the
elephants above mentioned from Mr Grote as the fore-
runner of the rainy season, though he mentions a form
under that name on the 9th of October. It was an
excessive hard gale, with violent thunder, lightning,
and rain; of which left there fell two inches in no
more than four hours. Neither is the quantity of
thunder and lightning at all comparable to what people
unacquainted with hot climates might be apt to ex-
pect. The only thunder-storms mentioned in the
journal were on May 31st, June 3d, 5th, 12th, 14th;
September 7th, October 9th, an elephant; and some
thunder on the 15th of the fame month.

The vegetable productions of Bombay are very insignificant. Mr Ives says, that its "soil is so barren as not to produce any one thing worth mentioning;" but afterwards informs us, that its "natural produce is the cocoa-nut tree, from which they extract a liquor called toddy. This is soft and mild when drunk immediately: but if it stands long, it gathers strength, and becomes very intoxicating; whence probably arose the term toddy-headed. For each tree a tax of 20s. a-year is paid to the company, which is appropriated towards maintaining the garrison and ships of war."

Mr Grofe gives an account somewhat different—
“The parts, or non-out grows, make the muff con-
fidential part of the health property, being started
wherever the futuation and foil is favorable to them.
Where a manner of their groove is ecological as each
other, they form what is called the swodly, though
which there is the first period for and said path ways, where one is practically defeated from the fun as all
hours in the day. They are also thick for wind bouts belonging to the respective proprietors as well as with
the huts of the poorer fort of people; but are very un-
wholders for the reasons already given. As to the

Vol. III. Part I.

cocoa-nut tree itself, not all the minute descriptions I Bombay.
have met with in many authors seem to me to come up

As to the rice fields, they differ in value, according to the fineness and quantity of rice they produce. The growth of this grain has a particularity not unworthy of notice, viz. that as it loves a watery soil, so to whatever height the water rises, wherever it is planted, the growth of the rice keeps measure with it, even to that of 12 and 14 feet; the summit always appearing above the surface of the water. It is also remarked, that the eating of new rice affects the eyes. The fact is certain, though the physical reason of it is unknown.

Here and there are interspersed some few \_brab\_
\_trees\_, or rather wild palm trees (the word \_brab\_ being
derived from \_brabas\_, which in the Portuguese signifies
wild.) They bear an insipid kind of fruit, about the
bigness of a common pear; but the chief profit from
them is the toddy, or liquor drawn from them by in-
cisions at the top, of which the arrack is reckoned bet-
ter than that produced by the cocoa-nut trees. They
are generally near the sea-side, as they delight most in
a sandy soil. It is on this tree that the \_toddy birds\_, so
called from their attachment to it, make their exqui-
sitely curious nests, wrought out of the thinnest reeds
and filaments of branches, with an inimitable mecha-
nism. The birds themselves are about the size of a
partridge, but are of no value either for plumage, song,
or the table.

“This island is a strong instance of the benefits of
a good government, and a numerous population, by
not a spot of it remaining uncultivated: so that tho’
it is far from producing sufficient for the consumption
of its inhabitants, and notwithstanding its many disad-
vantages of situation and soil, it yields incomparably
more than the adjacent island of Salette.”

Among the curiosities of Bombay Mr Ives mentions
a large terrapin or land tortoise kept at the governor's
house, the age of which was upwards of 200 years.
Frogs, which abound every where through the East
Indies, are very large at Bombay. Our author saw
one that measured 22 inches from the extremities of
the fore and hind feet when extended; and he suppo-
ses that its weight would not have been less than four
or five pounds. On the sea-shore round the island are
a great variety of beautiful shells, particularly the sort
called ventue-traps or evenue-traps, held in great es-
teem among the ladies some time ago. Several pounds
Sterling are said to have been given by a virtuoso for
one of these shells when the Commodore Lellie's col-
lection of shells was sold by auction.

Mr Ives enumerates the following kinds of snakes found on this island and other parts of the British em-pire in the East Indies. 1. The cobra de capella, growing from four to eight or nine feet long. They kill by their bite in 15 minutes. 2. The cobra manila is a small bluish snake, of the size of a man's little finger, and about a foot long, frequently seen about old walls. A species of these found at Bombay kill much sooner than


\section*{Source record}
\begin{description}
\item[Scan ID] \texttt{britannica:309f765b1ecc83e0a7f7762e44b4a927a7555443:train:2131}
\item[Image file] \texttt{page-2131-f30c24bfdf3b.jpg}
\item[Image SHA-256] \texttt{f30c24bfdf3b22773525a1262af33f8a656dd6dbb04d030b3d3e030cbde5405a}
\item[Dataset] NationalLibraryOfScotland/encyclopaedia\_britannica\_illustrated at revision \texttt{309f765b1ecc83e0a7f7762e44b4a927a7555443} (CC0-1.0)
\item[Extraction] model \texttt{nvidia/nemotron-3.5-parse}, raw reply \texttt{extract/raw/page-2131-f30c24bfdf3b.json}
\item[Raw reply SHA-256] \texttt{cf4c2f4d4ad5db69255ee7a04ceed96623f8687b8aeacf49ad5244bf1ddeb0d8}
\item[Correction] model \texttt{nvidia/Nemotron-3.5-Super-VL}, block crops, raw replies \texttt{correct/raw/crop\_nothink/page-2131-f30c24bfdf3b-b*.json}; super\_vl: 10
\item[Running header] BOM [ 353 BOM
\end{description}
\end{document}
